\section{One prime-splitting test for six Dedekind zeta functions}\label{sec:joint-test}

Fix $D=-12451$ or $-37363$. Let $K_1,\ldots,K_6$ be the six reflected
quintic fields for this $D$. At an unramified prime $p$, choose any two
distinct fields $K_i,K_j$ and record whether $p$ splits completely in
both. The same primes pass this test for every choice of the pair.
Together with a residue-class condition, this single bit gives a
one-sided criterion for RH in all six fields.

\paragraph{Joint Galois group.}
Choose virtual units $\alpha_1,\alpha_2$ representing a basis of
$\operatorname{Cl}(E)[5]$. Over $B=E(\zeta_5)$ their radical classes
are independent by restriction and norm. Adjoining both fifth roots
gives $M/B$ with translation group $V=\mathbb F_5^2$ by Kummer theory.
Choose the cyclotomic automorphisms fixing $E$ and both radicals,
and let $c$ be complex conjugation. As in the single-line case $c$ is
central. The six line evaluations of $V$ separate points, and each
quintic splitting field detects the full scalar group. Thus the joint
normal closure $T$ is $M^{\langle c\rangle}$ and
\[
 \operatorname{Gal}(T/\mathbb Q)
   = V\rtimes\mathbb F_5^\times,\qquad
 H=T^V=B^{\langle c\rangle}.
\]
The scalar group multiplies both coordinates. On the five embeddings
of the field associated with a radical line, a group element $(v,a)$
acts by $j\mapsto aj+\lambda(v)$, where $\lambda$ is a representative of the corresponding
line in the dual space. Its scalar at $p\nmid5D$ is
$\mu(p)=\chi_D(p)(p\bmod 5)\in\mathbb F_5^\times$.
Indeed Frobenius acts on $\zeta_5$ by raising to $p$ and on $E$ via
$\chi_D$. Simultaneous complex conjugation induces the identity on $H$.

For $a\ne1$ the affine map is conjugate to multiplication by $a$.
For $a=1$ a nonzero translation acts as a five-cycle. Of the six
functional lines, exactly one vanishes on a specified $v\ne0$.
Chebotarev gives the following natural densities; the second column
describes residue degrees.
\[
\begin{array}{lll}
 (v,a) & \text{splitting across the six fields} & \text{density}\\
 \hline
 v=0,\ a=1 & (1,1,1,1,1)\text{ in all six} & 1/100\\
 v\ne0,\ a=1 & (1,1,1,1,1)\text{ in one, }(5)\text{ in five}
      & 24/100\\
 a=-1 & (1,2,2)\text{ in all six} & 25/100\\
 a=2\text{ or }3 & (1,4)\text{ in all six} & 50/100
\end{array}
\]
Each of the six choices of the uniquely splitting quintic in the
second row occurs with density $4/100$. These densities hold
unconditionally for the two specified $D$.

\paragraph{The test.}
Let $S$ be a fixed finite set of excluded primes, containing those
dividing $5D$. When using reductions of defining polynomials one may
also exclude primes dividing their polynomial discriminants. Put
\[
 W(p)=25\,\mathbf1_{\{p\text{ splits completely in }K_i,K_j\}}
               -\mathbf1_{\{\mu(p)=1\}}, \qquad
 S_W(x)=\sum_{\substack{p\le x\\p\notin S}}
           W(p)\log p\,(p^{-1}-x^{-1}).
\]
The values of $W$ in the four rows are $24,-1,0,0$. There exists a
constant $c_W=\lim_{x\to\infty}S_W(x)$. Given Dirichlet GRH, the
following statements are equivalent.
\begin{enumerate}
\item All six $\zeta_{K_i}$ satisfy RH.
\item $\zeta_T$ satisfies RH.
\item For every $\epsilon>0$ there exists $C_\epsilon>0$ such that
$S_W(x)-c_W\le C_\epsilon x^{-1/2+\epsilon}$ for all sufficiently large $x$.
\end{enumerate}
Under these conditions, $|S_W(x)-c_W|\ll_{D,S} x^{-1/2}\log^2(2x)$.

\paragraph{Factorization and forward bound.}
Let $\tau_i$ be the character of the permutation representation on
the five embeddings of $K_i$ with one trivial character removed.
The table gives $W=\sum_i\tau_i$ on all Frobenius classes.
The regular character of $\operatorname{Gal}(T/\mathbb Q)$ equals the
inflated regular character of $\operatorname{Gal}(H/\mathbb Q)$ plus
$4\sum_i\tau_i$. Indeed their sum has value $100$ at the identity and
zero elsewhere. Artin formalism for permutation characters gives
\[
 \zeta_T(s)=\zeta_H(s)F(s)^4,\qquad
 F(s)=\prod_{i=1}^{6}\frac{\zeta_{K_i}(s)}{\zeta(s)}.
\]
Since $H/\mathbb Q$ is abelian, its zeta function factors into Dirichlet
$L$-functions. Dirichlet GRH makes each ratio in $F$ holomorphic on
$\Re s>1/2$, with a nonzero removable value at one, and gives the needed
zero-free region for $\zeta_H$ away from its pole. Zeros of the six
numerators in the right half of the strip therefore contribute positive
multiplicities to the product. The Dedekind functional equations give
the equivalence of the first two assertions.

Write $\theta_W(t)=\sum_{p\le t,\ p\notin S}W(p)\log p$. For each
$K_i$ the permutation trace counts degree-one prime ideals above $p$.
The fixed-field prime ideal theorem, applied to these six fields and
$\mathbb Q$, yields $\theta_W(t)=O(t\exp(-b\sqrt{\log t}))$ for some
$b>0$. The six main terms cancel. Higher residue degrees and prime
powers contribute $O_{D,S}(\sqrt t\log^2(2t))$, by summing over
$p^j\le t$, $j\ge2$. A possible real exceptional zero of a fixed field
is absorbed in the displayed asymptotic error. Under the first assertion,
the resulting bound is $\theta_W(t)=O_{D,S}(\sqrt t\log^2(2t))$.
Termwise integration gives
\[
 c_W=\int_1^\infty \theta_W(t)t^{-2}\,dt,\qquad
 R(x):=S_W(x)-c_W=-\int_x^\infty\theta_W(t)t^{-2}\,dt.
\]
This proves existence and the asserted two-sided estimate.
\href{https://arxiv.org/abs/1312.4463}{Greni\'e and Molteni,
\emph{Explicit versions of the prime ideal theorem for Dedekind zeta functions under GRH}}.

\paragraph{Converse from an upper bound.}
For $\Re s>1$ let $A(s)=\sum_{p\notin S}W(p)(\log p)p^{-s}$.
Euler factors give $A=-F'/F-J$ with $J$ holomorphic on $\Re s>1/2$.
At good primes the terms with powers $p^j$, $j\ge2$, are bounded by
$24(\log p)p^{-j\Re s}$ and converge locally uniformly.
The local factors at the finitely many excluded primes have inverse
roots of absolute value at most one, so their corrections are
holomorphic on $\Re s>0$. The prime ideal theorem gives $A(1)=c_W$
by partial summation. The Mellin transform of the centered sum is
\[
 P(s)=\int_1^\infty R(x)x^{-s}\,dx
     =\frac{A(s)}{s(s-1)}-\frac{c_W}{s-1}\qquad(\Re s>1).
\]
It is regular at one. For $\Re s>1/2$, any zero $\rho$ of $F$ of order
$m$ gives a simple pole of $P$ with residue $m/(\rho(1-\rho))$.
All its real poles in this region have positive residue, since a real
zero there lies below one.

Assume the third assertion and fix $1/2<\gamma<1$. Choose $C,X$
so $U(x)=Cx^{\gamma-1}-R(x)\ge0$ for $x\ge X$.
The Mellin integral $\int_X^\infty U(x)x^{-s}\,dx$ agrees with
$C/(s-\gamma)-P(s)$ up to an entire function, and its abscissa of
convergence is at most one. If that abscissa exceeded $\gamma$,
Landau's lemma for a nonnegative Mellin integrand would force a
singularity at the real abscissa. The meromorphic continuation permits
only a simple pole of negative residue there. Approaching the pole
from the right would make the integral negative. This contradiction
puts the abscissa at most at $\gamma$, so the integral is holomorphic
throughout $\Re s>\gamma$, where $P$ has no poles. Taking all
$\gamma>1/2$ proves zero-freeness of $F$ there. This proves the
converse. The prime estimate itself remains to be proved for these
nonabelian fields.

\paragraph{The information missing from residue classes.}
For every positive integer $m$ and every complex-valued function $h$
on $(\mathbb Z/m\mathbb Z)^\times$, Chebotarev gives
\[
 \lim_{x\to\infty}\frac1{\pi(x)}
 \sum_{\substack{p\le x\\p\notin S,\ p\nmid m}}
       |W(p)-h(p\bmod m)|^2
  =6+\frac1{\varphi(m)}
           \sum_{a\in(\mathbb Z/m\mathbb Z)^\times}|h(a)|^2 .
\]
Indeed $[V\rtimes\mathbb F_5^\times,V\rtimes\mathbb F_5^\times]=V$,
since the commutators with scalar $2$ give every translation. Thus
$T\cap\mathbb Q(\zeta_m)$ is contained in $H$. Conditional on a
cyclotomic Frobenius, the compatible elements of $\operatorname{Gal}(T/\mathbb Q)$
are uniformly counted on a union of full translation cosets. The mean
of $W$ on each coset is zero: on $V$ its sum is $24-24=0$, and on
the other three cosets $W$ vanishes. The mixed term in the square
therefore has zero mean. The first two rows of the table give the
unconditional mean square $(24^2+24)/100=6$.

The cubic case appears in
\href{https://github.com/VladimirReshetnikov/ProveIt/tree/1e1728ad221d69e48fea09bf754b73422af44943/docs/incoming}
{the incoming cubic-field report}. The Mellin argument above uses
\href{https://github.com/lsernau/grh-lean/blob/72fead102f663d5c7a407964a13dc0cf9d977207/grh-compact/src/New/Endgame/OneSidedCriterion.lean}
{the Lean repository's Landau criterion}.
